\section{Introduction}
\label{sec:intro}

Enterprise collections contain regulations, failure analyses and operating
procedures. A missed answer can originate before passage ranking: a vague
catalogue description sends the agent to the wrong collection, or a stale
index points to a document's former path. The operator needs to tell these
failures apart because they require different repairs.

Document acquisition, organisation and retrieval are often studied
separately. Enterprise content management studies document
governance \cite{paivarinta-ecm-2005,tyrvaainen-ecm-2006}, while metadata
quality and tagging affect how users discover documents
\cite{park-metadata-quality-2009,golder-tagging-2006}. Distributed
information retrieval addresses the selection of collections
\cite{callan1995searching,si2003redde}. Retrieval-augmented generation
(RAG) uses ranked passages to inform generation, and modular retrieval
workflows can expose evidence through interfaces such as the Model
Context Protocol (MCP) \cite{rag-lewis,dpr,modular-rag,mcp,gao-rag-survey-2023}.
This paper asks how management state can make retrieval decisions
inspectable: whether a document is visible to search, whether an agent
reaches it and which source text the workflow returns.

The Hierarchical Management Retrieval System (HMRS) is a self-hosted
platform that exposes the collection membership and descriptions that
operators maintain, rather than deriving a second hierarchy at query time.
Agents use this state to navigate towards source text or search within a
specified collection. Both paths expose source references so that
operators can inspect the evidence rather than infer it from the answer.

\noindent\textbf{Contributions.}
\begin{enumerate}
\item \textbf{A management model for evidence access.} HMRS connects
three managed levels (collection, document and segment) to four access
views (catalogue, shelf, document and segment). Collection membership
specifies the intended scope for scoped calls; descriptions guide selection
and source references link passages to documents (Section~\ref{sec:system}).
\item \textbf{A content-informed document lifecycle with post-operation
probes.} The acquisition workflow derives placement, descriptions and
tags from parsed content. Operators can reconcile files, tree entries
and library metadata after lifecycle changes and probe retrieval for
stale index entries
(Sections~\ref{sec:acquisition}--\ref{sec:organization}).
\item \textbf{Agent access policies over shared state.} MCP tools let an
agent navigate from collection and document descriptions to source text,
or search the indices of specified collections directly
(Section~\ref{sec:retrieval}).
\item \textbf{An evidence gate that keeps admission and ordering as
separate, inspectable decisions.} An optional server-side gate fetches
candidate text, scores passages and returns those that meet the threshold
with source references. It reports fetch and scoring failures separately
so that the caller can report an incomplete review rather than an
unsupported not-found result (Section~\ref{sec:gate}). Candidate order is
controlled separately: in the replay, using verifier scores to rank
candidates displaced gold documents from the top five
(Section~\ref{sec:e6}).
\end{enumerate}

The evaluation tests what each checkpoint reveals. Index probes
separate index visibility from the result seen by the caller. Routing and containment
experiments separate reaching the right collection from delivering the
right text. A fixed-pool replay tests document ordering without changing
candidate membership. The native answer comparison supplies deployment
context; it does not substitute for the catalogue-first tests.

\subsection{Research questions and reading guide}
\label{sec:rq}

\looseness=-1 \textbf{RQ1} examines document acquisition, collection membership and
lifecycle consistency. \textbf{RQ2} measures collection selection under
catalogue-first access. \textbf{RQ3} compares target-string containment
and retrieval text volume across the tested evidence configurations.
\textbf{RQ4} asks how verifier-score ordering changes gold-document
Hit@5 in fixed candidate pools. \textbf{RQ5} examines abstention and
operational failures. \textbf{RQ6} compares judged answer quality under
native configurations.

Section~\ref{sec:related} positions HMRS within related work.
Sections~\ref{sec:system}--\ref{sec:implementation} describe the management
model and its implementation. Section~\ref{sec:eval} reports the evaluation,
followed by the discussion, limitations and conclusions in
Sections~\ref{sec:discussion}--\ref{sec:conclusion}.
